% ========================================================================
% فصل ۲: شبکه اقتضایی خودرویی چیست؟ (بدنهٔ فصل، خروجی گره sam-work/sections)
% ========================================================================
\section{شبکه اقتضایی خودرویی چیست؟}

\subsection{تعریف و مفهوم}

شبکه اقتضایی خودرویی زیرمجموعه‌ای از شبکه‌های اقتضایی سیار \LTRfootnote{\lr{Mobile Ad-hoc Networks (MANET)}} است که ارتباط میان خودروهای نزدیک به هم و نیز ارتباط میان خودروها و زیرساخت‌ها را تسهیل می‌کند \cite{hozouriOverviewVANETVehicular2023}. آنچه این شبکه را از دیگر اعضای خانواده شبکه‌های اقتضایی، مانند شبکه‌های اقتضایی پرنده \LTRfootnote{\lr{Flying Ad-hoc Networks (FANET)}} و شبکه‌های حسگر بی‌سیم \LTRfootnote{\lr{Wireless Sensor Networks (WSN)}}، متمایز می‌کند، تحرک بسیار زیاد گره‌ها، وابستگی به زیرساخت حمل‌ونقل، جابه‌جایی پویای شبکه و اتصال پراکنده و منقطع است \cite{hozouriOverviewVANETVehicular2023}. در نگاه سنتی، هر خودرو گره‌ای شبکه به شمار می‌آید که داده‌ها را از طریق ارتباط خودرو به خودرو و خودرو به زیرساخت ارسال یا بازپخش می‌کند؛ ارتباط خودرو به خودرو به خودروها امکان می‌دهد اطلاعات ایمنی را مستقیماً با همسایگان خود به اشتراک بگذارند و ارتباط خودرو به زیرساخت برای گردآوری اطلاعات از تجهیزات کنار جاده به کار می‌رود \cite{xuInternetVehiclesBig2018}.

این تصویر سنتی در سال‌های اخیر تغییر کرده است. خودروهای امروزی به واحدهای پردازشی توانمند و حافظه‌های بزرگ مجهز شده‌اند و دیگر فقط یک گره متحرک نیستند، بلکه خود به مرکزی برای محاسبه و ذخیره‌سازی با توان پردازش هوشمند تبدیل شده‌اند. با گسترش چشمگیر مقیاس شبکه و نیاز به پردازش هم‌زمان بلادرنگ و بلندمدت داده‌ها، شبکه‌های اقتضایی خودرویی سنتی به سمت اینترنت خودروها \LTRfootnote{\lr{Internet of Vehicles (IoV)}} در حال تحول‌اند \cite{xuInternetVehiclesBig2018}. در همین چارچوب، خودروی متصل به خودرویی گفته می‌شود که با امکان ارتباط بی‌سیم با محیط داخلی و خارجی خود تعامل دارد؛ یعنی هم با حسگرهای درون خودرو، هم با خودروهای دیگر، هم با زیرساخت جاده‌ای و هم با اینترنت در ارتباط است. این تعامل‌ها آگاهی خودرو از وضعیت پیرامون را افزایش می‌دهند و محیطی سرشار از اطلاعات برای راننده و سرنشینان فراهم می‌کنند؛ خودروهای متصل را سنگ‌بنای اینترنت خودروها و نسل آینده سامانه‌های حمل‌ونقل هوشمند \LTRfootnote{\lr{Intelligent Transportation Systems (ITS)}} می‌دانند \cite{luConnectedVehiclesSolutions2014}. اهمیت اقتصادی این حوزه نیز از مدت‌ها پیش آشکار بوده است؛ برای نمونه، در سال ۲۰۱۴ و به استناد یک گزارش تجاری پیش‌بینی شده بود که بازار جهانی خودروهای متصل تا سال ۲۰۱۹ به ۱۳۱٫۹ میلیارد دلار برسد \cite{luConnectedVehiclesSolutions2014}.

\subsection{ارکان اصلی شبکه اقتضایی خودرویی}

تعریف مفهوم، طبیعتاً پرسش اجزای تشکیل‌دهنده شبکه را مطرح می‌کند. در مدل پایه‌ای که بیشتر طرح‌های امنیتی و حریم خصوصی بر آن بنا شده‌اند، شبکه اقتضایی خودرویی از سه موجودیت اصلی تشکیل شده است: واحد سوارشونده در خودرو، واحد کنارجاده‌ای و مرجع قابل اعتماد \cite{rahmawatiagustinaSystematicLiteratureReview2025}.

واحد سوارشونده \LTRfootnote{\lr{On-Board Unit (OBU)}} سخت‌افزاری است که در هر خودرو نصب می‌شود و ارتباط آن خودرو را با واحدهای کنارجاده‌ای و واحدهای سوارشونده سایر خودروها فراهم می‌کند. این واحد پیام‌ها را از منبعی که می‌تواند خودروی دیگر یا حسگری باشد دریافت و اعتبارسنجی می‌کند و سپس از طریق سامانه ارتباطات اختصاصی کوتاه‌برد \LTRfootnote{\lr{Dedicated Short-Range Communications (DSRC)}} برای خودروهای دیگر پخش می‌کند \cite{hozouriOverviewVANETVehicular2023}. واحد سوارشونده تنها بخشی از تجهیزات یک خودروی هوشمند است؛ چنین خودرویی معمولاً پردازنده مرکزی، فرستنده‌گیرنده بی‌سیم، گیرنده موقعیت‌یاب جهانی \LTRfootnote{\lr{Global Positioning System (GPS)}}، حسگرهایی مانند حسگر فراصوت، رادار، لیدار \LTRfootnote{\lr{LIDAR}} و دوربین پایش راننده نیز دارد و وظیفهٔ واحد سوارشونده تنظیم ارتباط خودرو با واحدهای کنارجاده‌ای و سایر خودروها از طریق ارتباطات اختصاصی کوتاه‌برد و استانداردهای سلولی است \cite{hozouriOverviewVANETVehicular2023}.

واحدهای کنارجاده‌ای \LTRfootnote{\lr{Road-Side Unit (RSU)}} واحدهای ثابتی‌اند که در کنار جاده‌ها نصب می‌شوند و واحدهای سوارشونده از طریق ارتباط خودرو به زیرساخت با آن‌ها تعامل می‌کنند \cite{hozouriOverviewVANETVehicular2023}. این واحدها اطلاع‌رسانی‌ها و دسترسی به اینترنت را در اختیار خودروها می‌گذارند و اطلاعات خودروهای حاضر در محدوده پوشش خود را گردآوری کرده و برای تحلیل به ابر می‌فرستند \cite{hozouriOverviewVANETVehicular2023}. ارتباط با زیرساخت جاده‌ای برای جلوگیری از تصادفات یا کاهش پیامدهای آن‌ها و نیز برای مدیریت کارآمد سامانه‌های حمل‌ونقل هوشمند حیاتی است؛ فناوری ارتباطات اختصاصی کوتاه‌برد در این زمینه ابزار اصلی ارتباط خودرو با چراغ‌های راهنمایی، تابلوها و حسگرهای کنار جاده است. با این حال، واحد کنارجاده‌ای لزوماً نقش دروازه اینترنت را ندارد و می‌تواند صرفاً ارائه‌دهنده محتوا یا خدمات محلی باشد \cite{luConnectedVehiclesSolutions2014}.

مرجع قابل اعتماد \LTRfootnote{\lr{Trusted Authority (TA)}} نهادی مرکزی است که مسئولیت مدیریت و امنیت ارتباطات را در شبکه بر عهده دارد. این مرجع با صدور اعتبارنامه‌های امنیتی، مانند گواهی‌های دیجیتال و شبه‌نام‌ها، تضمین می‌کند که فقط موجودیت‌های مجاز در شبکه مشارکت کنند و در صورت لزوم به نهادهای قانونی در ردیابی خودروهای مخرب کمک می‌کند. همین جایگاه مرکزی نقطه ضعف آن نیز هست: از دسترس خارج شدن مرجع قابل اعتماد بر همه موجودیت‌های شبکه اثر می‌گذارد و از آنجا که این مرجع هویت واقعی همه موجودیت‌ها را می‌داند، به هدفی جذاب برای مهاجمان تبدیل می‌شود و نفوذ به آن می‌تواند به نقض گسترده حریم خصوصی بینجامد. به همین دلیل، در برخی طرح‌ها وظایف این مرجع میان دو نهاد جداگانه، یعنی مرجع ردیابی \LTRfootnote{\lr{Trace Authority (TRA)}} و مرکز تولید کلید \LTRfootnote{\lr{Key Generation Center (KGC)}}، تقسیم شده است \cite{rahmawatiagustinaSystematicLiteratureReview2025}. این ضعف ساختاری، چنان‌که در فصل‌های بعد خواهیم دید، یکی از انگیزه‌های اصلی روی آوردن به راهکارهای غیرمتمرکز است.

\subsection{معماری و انواع ارتباطات}

اجزای یادشده در قالب معماری شبکه به یکدیگر وصل می‌شوند؛ شکل \ref{fig:vanet_arch} معماری کلی شبکه اقتضایی خودرویی و انواع ارتباطات خودرو با هر موجودیت \LTRfootnote{\lr{Vehicle-to-Everything (V2X)}} را نشان می‌دهد.

\begin{figure}[H]
\centering
\includegraphics[width=0.9\textwidth]{VANET artitecture.png}
\caption{معماری شبکه اقتضایی خودرویی و انواع ارتباطات خودرو با هر موجودیت}
\label{fig:vanet_arch}
\end{figure}

پایه‌ای‌ترین نوع این ارتباطات، ارتباط خودرو به خودرو \LTRfootnote{\lr{Vehicle-to-Vehicle (V2V)}} است که بدون نیاز به هیچ زیرساختی و مستقیماً میان خودروها برقرار می‌شود و هدف اصلی آن تبادل اطلاعات مرتبط با ایمنی است. در ارتباط خودرو به زیرساخت \LTRfootnote{\lr{Vehicle-to-Infrastructure (V2I)}}، خودرو با واحدهای کنارجاده‌ای یا ایستگاه‌های پایه شبکه سلولی ارتباط می‌گیرد تا اطلاع‌رسانی‌ها را دریافت کند و به اینترنت دسترسی یابد. ارتباط زیرساخت به زیرساخت \LTRfootnote{\lr{Infrastructure-to-Infrastructure (I2I)}} برای تبادل داده‌های ترافیکی میان تجهیزات جاده‌ای به کار می‌رود. ارتباط خودرو به عابر پیاده \LTRfootnote{\lr{Vehicle-to-Pedestrian (V2P)}} از طریق دستگاه‌های قابل‌حمل افراد برقرار می‌شود و ارتباط خودرو به موانع کنار جاده \LTRfootnote{\lr{Vehicle-to-Barrier (V2B)}} نیز از همین دسته‌اند. ارتباط خودرو به ابر \LTRfootnote{\lr{Vehicle-to-Cloud (V2C)}} واحدهای کنارجاده‌ای را به سرورهای ابری وصل می‌کند تا پردازش داده، تصمیم‌گیری و پیش‌بینی وضعیت حمل‌ونقل در آنجا انجام شود. ارتباط خودرو به پهپاد \LTRfootnote{\lr{Vehicle-to-UAV (V2U)}} پیوندی میان زمین و هوا ایجاد می‌کند و ارتباط خودرو به حسگر \LTRfootnote{\lr{Vehicle-to-Sensor (V2S)}} خودرو را به حسگرهای تعبیه‌شده وصل می‌کند \cite{hozouriOverviewVANETVehicular2023}.

نام‌گذاری این ارتباطات در همه منابع یکسان نیست؛ برای مثال، در برخی منابع \lr{V2I} به معنای ارتباط خودرو با اینترنت به کار رفته و ارتباط با زیرساخت جاده‌ای \lr{V2R} نامیده شده است \cite{luConnectedVehiclesSolutions2014}. نیازهای کیفیت خدمت \LTRfootnote{\lr{Quality of Service (QoS)}} در این ارتباطات نیز یکسان نیست: کاربردهای غیرایمنی به گذردهی بالا نیاز دارند، در حالی که خدمات ایمنی باید با تأخیر کم و قابلیت اطمینان بالا در دسترس باشند \cite{hozouriOverviewVANETVehicular2023}.

\subsection{اهمیت شبکه اقتضایی خودرویی}

با شناخت معماری و انواع ارتباطات، به کاربردها و خدماتی می‌رسیم که اهمیت این شبکه را روشن می‌سازد. شبکه اقتضایی خودرویی برای ارائه هشدارهای ایمنی و خدمات اطلاعاتی و سرگرمی کم‌تأخیر به رانندگان و خودروهای بدون راننده طراحی شده است و کاربردهای آن را می‌توان در چهار دسته ایمنی، اطلاعات و سرگرمی، بهبود ترافیک و پایش رانندگی جای داد \cite{hozouriOverviewVANETVehicular2023}.

در این میان، ایمنی ترافیک مهم‌ترین انگیزه توسعه این شبکه‌هاست. کاربردهای ایمنی شامل کمک به راننده، رساندن اطلاعات ایمنی و هشدار به راننده‌اند و نمونه‌هایی چون جلوگیری از برخورد، هشدار برخورد در تقاطع، ترمز اضطراری، هشدار تغییر خط، هشدار بسته بودن جاده، هشدار عبور عابر پیاده و اطلاع‌رسانی پس از تصادف را در بر می‌گیرند. پیام‌های هشدار می‌توانند وضعیت و زمان‌بندی چراغ راهنمایی، حجم ترافیک، حق تقدم عبور و شرایط آب‌وهوایی را نیز منتقل کنند \cite{hozouriOverviewVANETVehicular2023}. پیام ایمنی معمولاً حاوی موقعیت، سرعت، جهت حرکت و شتاب خودرو است، اندازه‌ای در حدود ۲۰۰ تا ۵۰۰ بایت دارد و باید با محدودیت‌های سخت‌گیرانه تأخیر و قابلیت اطمینان تحویل داده شود \cite{xuInternetVehiclesBig2018}.

دسته دوم، بهبود جریان ترافیک است. مدیریت تقاطع‌ها، مدیریت ازدحام و اطلاع‌رسانی درباره وضعیت جاده از کاربردهای این دسته به شمار می‌روند \cite{hozouriOverviewVANETVehicular2023}. خودروها می‌توانند از طریق ارتباط با زیرساخت سامانه حمل‌ونقل هوشمند وضعیت ترافیک را پیش‌بینی کنند و بهترین مسیر را بیابند، و اطلاعات ترافیکی که از داده‌های حرکت خودروها استخراج می‌شود، امکان ناوبری بلادرنگ را فراهم می‌کند. اطلاع‌رسانی درباره وضعیت چراغ‌ها، ترافیک مسیر پیش رو و پارکینگ‌های خالی اطراف نیز به بهبود تجربه رانندگی و کارایی جاده کمک می‌کند \cite{xuInternetVehiclesBig2018}. اهمیت این موضوع با یک عدد روشن‌تر می‌شود: هزینه زمان سفر و سوخت اضافی ناشی از ازدحام در ۴۹۸ منطقه شهری آمریکا در سال ۲۰۱۱ به ۱۲۱ میلیارد دلار رسیده بود، در حالی که این رقم در سال ۱۹۸۲ حدود ۲۴ میلیارد دلار بود \cite{luConnectedVehiclesSolutions2014}.

دسته سوم، خدمات اطلاعاتی و سرگرمی است. کاربردهایی مانند تماس صوتی اینترنتی، بازی برخط، اشتراک ویدئو، دسترسی به اینترنت، به‌روزرسانی خودکار نقشه‌ها و اطلاعات گردشگری در این دسته قرار می‌گیرند؛ این کاربردها تأخیر را تا حدی تحمل می‌کنند اما ممکن است پهنای باند زیادی بخواهند \cite{hozouriOverviewVANETVehicular2023}. خدمات وابسته به مکان، مانند معرفی نقاط دیدنی و بهینه‌سازی مسیر، نیز بخشی از همین دسته‌اند. پشتوانه این تقاضا روشن است: مردم انتظار دارند در خودروی خود همان اتصالی را داشته باشند که در خانه و محل کار دارند \cite{luConnectedVehiclesSolutions2014}.

سرانجام، این شبکه‌ها نقش مهمی در تحقق خودروهای خودران دارند. داده‌های گسترده‌ای که در اینترنت خودروها جریان دارد، فناوری کلیدی برای خودروهای خودران است، زیرا این خودروها داده‌های حسگرهای خود مانند دوربین، رادار، لیدار و موقعیت‌یاب را با اطلاعات دریافتی از خودروهای متصل دیگر، از جمله وضعیت جاده و ترافیک، ترکیب می‌کنند. پیش‌بینی شده است که یک خودروی خودران بیش از یک ترابایت داده در ساعت تولید کند، و امنیت و حریم خصوصی این داده‌ها یکی از نگرانی‌های جدی در این حوزه است \cite{xuInternetVehiclesBig2018}. همین نگرانی پلی است به موضوع اصلی این گزارش، یعنی امنیت و حریم خصوصی در شبکه‌های خودرویی.

\subsection{ویژگی‌های متمایز}

کاربردهای گسترده این شبکه، از ویژگی‌های ذاتی آن سرچشمه می‌گیرد. شبکه اقتضایی خودرویی ویژگی‌هایی دارد که آن را از دیگر شبکه‌های بی‌سیم متمایز می‌کند.

نخستین و آشکارترین ویژگی، تحرک بالای گره‌هاست. البته تحرک در این شبکه یکنواخت نیست: واحدهای کنارجاده‌ای ثابت‌اند، خودروها در ترافیک شهری و تقاطع‌ها آهسته حرکت می‌کنند و در بزرگراه‌ها سرعت زیادی دارند \cite{hozouriOverviewVANETVehicular2023}. در عین حال، حرکت خودروها به شبکه جاده‌ها محدود است و همین محدودیت باعث می‌شود که حرکت خودرو در بازه زمانی مشخصی تا حدی قابل پیش‌بینی باشد \cite{hozouriOverviewVANETVehicular2023,luConnectedVehiclesSolutions2014}.

پیامد مستقیم تحرک زیاد، تغییرات مکرر و سریع توپولوژی شبکه است؛ هر خودرو مسیر حرکت خاص خود را دارد و ساختار شبکه پیوسته در حال تغییر است \cite{luConnectedVehiclesSolutions2014}. پویایی توپولوژی در کنار برد محدود ارتباط خودرو به خودرو می‌تواند به تکه‌تکه شدن مکرر شبکه و قطع جریان داده بینجامد، و موانعی مانند ساختمان‌ها و کامیون‌ها نیز اتصال را متناوب می‌کنند \cite{luConnectedVehiclesSolutions2014}. کاهش چگالی خودروها هم شبکه را به بخش‌های جدا از هم تقسیم می‌کند و بر تحویل بسته‌ها اثر می‌گذارد \cite{hozouriOverviewVANETVehicular2023}. این مسائل در فصل بعد، ذیل چالش‌های شبکه‌سازی، با جزئیات بیشتری بررسی می‌شوند.

ویژگی سوم به منابع انرژی و پردازشی مربوط است. در شبکه‌های حسگر بی‌سیم، تأمین انرژی گره‌ها چالش اصلی است \cite{hozouriOverviewVANETVehicular2023}، اما در شبکه اقتضایی خودرویی ارتباطات با محدودیت انرژی روبه‌رو نیست و هر خودرو می‌تواند توان پردازشی نسبتاً بالایی داشته باشد \cite{luConnectedVehiclesSolutions2014}. واحدهای سوارشونده انرژی خود را به‌طور پیوسته از باتری خودرو می‌گیرند و به همین دلیل ارتباط میان گره‌ها محدودیت جدی از نظر توان و حافظه ندارد \cite{hozouriOverviewVANETVehicular2023}. با این حال، از منظر طراحی راهکارهای امنیتی، توان محاسباتی خودرو در مقایسه با رایانه، تبلت و تلفن هوشمند محدودتر است و برخی سازوکارهای امنیتی به دلیل سربار بالا قابل اجرا نیستند \cite{farsimadanReviewSecurityChallenges2025}؛ بنابراین محدودیت منابع پردازشی ویژگی‌ای نسبی است که در طراحی رمزنگاری سبک باید در نظر گرفته شود.

ویژگی چهارم، نیاز به ارتباط بلادرنگ و کم‌تأخیر است. خدمات ایمنی باید با تأخیر کم و قابلیت اطمینان بالا در دسترس باشند \cite{hozouriOverviewVANETVehicular2023,luConnectedVehiclesSolutions2014} و به همین دلیل تحمل‌پذیری در برابر خطا اهمیت زیادی دارد؛ هر خطا در ارتباط می‌تواند به تأخیر در رسیدن هشدار و در نهایت به تصادف منجر شود \cite{hozouriOverviewVANETVehicular2023}.

در کنار این ویژگی‌ها، ناهمگونی گره‌ها و خدمات، گستره جغرافیایی تقریباً نامحدود شبکه، کمبود طیف فرکانسی، اثر محیط بر انتشار سیگنال و دقت محدود اطلاعات موقعیت نیز از ویژگی‌های این شبکه به شمار می‌روند؛ برای مثال، دقت سیگنال موقعیت‌یاب جهانی در شرایط مطلوب معمولاً بین ۵ تا ۳۰ متر است \cite{hozouriOverviewVANETVehicular2023}.

دو ویژگی آخر، امنیت و حریم خصوصی داده‌هاست که موضوع اصلی این گزارش است. داده‌هایی که در شبکه جابه‌جا می‌شوند باید رمزنگاری شوند، یکپارچگی بسته‌ها حفظ شود، اصالت فرستنده احراز شود و فرستنده نتواند ارسال پیام را انکار کند. با این حال، تغییر مداوم توپولوژی مدیریت کلیدها را بسیار دشوار می‌کند و ابطال کلیدها به فهرست‌های طولانی می‌انجامد. از سوی دیگر، کاربران تمایلی ندارند که اطلاعات خودرو و مقصدشان در اختیار دیگران قرار گیرد و طراح شبکه باید میان منافع عمومی و حریم خصوصی کاربران تعادل برقرار کند \cite{hozouriOverviewVANETVehicular2023}. فصل‌های بعدی این گزارش به همین دو مسئله و راهکارهای پیشنهادی برای آن‌ها اختصاص دارد.
